%!TEX program = xelatex
% =============================================================
%  北京工业大学学位论文模板 (BJUT Thesis Template)
%  使用方法：在下方 \bjutset{...} 中填写自己的论文信息，
%           然后在各 chapt*.tex / cabstract.tex / eabstract.tex
%           等文件中填写内容，最后用 XeLaTeX 编译。
%  编译流程： xelatex main → bibtex main → xelatex main → xelatex main
% =============================================================

\documentclass{bjutthesis}

% ---- 常用宏包（按需保留 / 增删）----
\usepackage{slashed}
\usepackage{bm,mathrsfs}
\usepackage{empheq}
\usepackage{microtype}
\usepackage{longtable,array}
% 索引支持（可选；博士论文常用，硕士一般不需要）
% 如不需要索引，注释掉下面这行及文末的 \printindex
\usepackage{makeidx}
\makeindex

\allowdisplaybreaks[2]

% ---- cleveref 引用名（中文） ----
\crefformat{equation}{式~(#2#1#3)}
\Crefformat{equation}{式~(#2#1#3)}
\crefname{table}{表}{表}
\Crefname{table}{表}{表}
\crefname{figure}{图}{图}
\Crefname{figure}{图}{图}
\crefformat{chapter}{第#2#1#3章}
\Crefformat{chapter}{第#2#1#3章}
\crefname{section}{节}{节}
\Crefname{section}{节}{节}
\crefname{theorem}{定理}{定理}
\Crefname{theorem}{定理}{定理}
\crefname{lemma}{引理}{引理}
\Crefname{lemma}{引理}{引理}

% ---- 定理类环境（中文） ----
\newtheorem{theorem}{定理}[chapter]
\newtheorem{lemma}[theorem]{引理}
\newtheorem{proposition}[theorem]{命题}
\newtheorem{corollary}[theorem]{推论}
\newtheorem{definition}[theorem]{定义}
\newtheorem{remark}[theorem]{注}
\newenvironment{proof}{\par\noindent\textit{证明.}\ }{\hfill$\square$\par}

% ---- 自定义命令（示例，按需修改）----
\newcommand{\dd}{\mathop{}\!\mathrm{d}}
\newcommand{\Var}{\operatorname{Var}}
\newcommand{\Cov}{\operatorname{Cov}}
\def\cappendixname{附录}

% 中文圆括号公式编号
\makeatletter
\renewcommand{\tagform@}[1]{\maketag@@@{（#1）}}
\makeatother

% =============================================================
%  论文信息（请填写）
% =============================================================
\bjutset{
  % ── 封面类型（决定封面底色）─────────────────────
  % 可选值（四选一，删除不需要的注释）：
  %   academic-master      学术学位硕士 / 同等学力硕士（蓝色）
  %   professional-master  全日制/非全日制专业学位硕士（青绿色）
  %   academic-doctor      学术学位博士（深红色）
  %   professional-doctor  专业学位博士（橙色）
  covertype      = {academic-master },         % ← 按实际修改
  % ── 论文题目 ─────────────────────────────────────
  ctitle         = {论文中文题目},            % 中文题目（≤ 25 字，不设副标题）
  etitle         = {English Title of the Thesis},  % 英文题目（字母全部大写）
  % ── 作者信息 ─────────────────────────────────────
  cauthor        = {作者姓名},                % 姓名（中文）
  studentid      = {S202200001},             % 学号（按实际填写）
  cdegree        = {理学碩士},              % 学位，按实际选填：
                                             %   理学硕士 / 工学硕士 / 管理学硕士
                                             %   理学博士 / 工学博士 / 管理学博士
  cdiscipline    = {学科名称},               % 一级学科，如：物理学 / 机械工程
  cmajor         = {专业方向},               % 研究方向 / 专业方向
  ccollege       = {所在学院},               % 学院名称，如：理学院 / 机械与精仪工程学院
  % ── 导师信息 ─────────────────────────────────────
  csupervisor    = {导师姓名},               % 导师姓名
  csupervstitle  = {教授},                   % 导师职称，如：教授 / 副教授（可留空 {}）
  % ── 时间与学校 ───────────────────────────────────
  cdate          = {2025年6月},              % 答辩 / 提交年月
  corganization  = {北京工业大学},            % 一般不改
  % ── 封面行政信息 ─────────────────────────────────
  schoolcode     = {10005},                  % 学校代码（北工大固定为 10005，勿改）
  secretlevel    = {公开},                   % 密级：公开 / 限制 / 秘密 / 机密
  clc            = {O413},                   % 中图分类号（在国家图书馆官网查询）
  udc            = {530},               % UDC 分类号（与中图对应，不确定可留空 {}）
}

\begin{document}

% ---- 图表标题中文设置 ----
\captionsetup[figure]{name=图}
\captionsetup[table]{name=表}
\renewcommand{\cftfigpresnum}{图~}
\renewcommand{\cfttabpresnum}{表~}

% ---- 封面与扉页 ----
\makecover
\cleardoublepage
\maketitle
\makestate

% ---- 前置部分：摘要、目录、符号表 ----
\frontmatter
\begin{cabstract}

\linespread{1.08}\selectfont

% ====== 在此填写中文摘要正文 ======
% 一般包括：研究背景与意义、研究内容与方法、主要结果与结论。

% ===================================

\ckeywords{关键词一；关键词二；关键词三；关键词四；关键词五}

\end{cabstract}
     % 中文摘要
\begin{eabstract}
\linespread{0.9}\selectfont

% ====== Write the English abstract here ======
% Content should correspond to the Chinese abstract.
% Target: approximately 300 content words (实词).
% Do not use formulas, tables, or citation numbers in the abstract.

% =============================================

% 3–5 keywords, separated by commas, ordered from broader to narrower scope.
\ekeywords{keyword one, keyword two, keyword three, keyword four, keyword five}
\end{eabstract}
     % 英文摘要
\makebitoc              % 中英文双语目录
% 符号表（可选：不需要则注释掉下一行）
\chapter*{物理量名称及符号表}
\addcontentsline{toc}{chapter}{物理量名称及符号表}

\begin{longtable}{>{\raggedright\arraybackslash}p{0.18\textwidth} >{\raggedright\arraybackslash}p{0.74\textwidth}}
\toprule
符号 & 含义 \\
\midrule
\endfirsthead
\toprule
符号 & 含义 \\
\midrule
\endhead
% 在此添加符号与含义条目。例如：
% $A$ & 符号 A 的含义。\\
\bottomrule
\end{longtable}
      % 物理量名称及符号表

% ---- 主体部分：正文章节 ----
\mainmatter
% 章节标题使用中文（默认即为「第 X 章」格式，由 ctexbook 提供）
\ctexset{
  chapter = {
    name      = {第,章},
    number    = {\arabic{chapter}},
    aftername = \space,
  }
}
\chapter{绪论}

\section{研究背景与意义}

% 在此描述研究背景与意义。示例段落（替换为自己的内容）：
在……领域中，……问题具有重要的理论与应用价值。
随着……的发展，传统方法已难以满足……的需求，亟需新的理论框架。

\section{国内外研究现状}

% 建议按研究方向分小节综述，每方向引用 3～5 篇代表性文献。
% 文献引用示例（需在 reference.bib 中有对应条目）：
Doe 等人~\cite{doe2020} 提出了……方法；
Smith~\cite{smith2021} 进一步将其推广至……情形。
目前仍存在以下不足：
（1）……；
（2）……；
（3）……。

\section{主要研究内容与论文结构}

本文各章节安排如下：

\begin{enumerate}
  \item \textbf{第一章}：绪论，介绍研究背景、现状与论文结构。
  \item \textbf{第二章}：……（本章概述）。
  \item \textbf{第三章}：……（本章概述）。
  \item \textbf{第四章}：……（本章概述）。
  \item \textbf{第五章}：……（本章概述）。
  \item \textbf{第六章}：全文总结与展望。
\end{enumerate}

% ============================================================
%  第二章：LaTeX 常用元素示例
%  本章展示公式、定理、表格、图片、交叉引用等常见用法。
%  实际写作时请将本章内容替换为自己的章节。
% ============================================================
\chapter{第二章标题}

% ────────────────────────────────────────────────────────────
\section{数学公式}
% ────────────────────────────────────────────────────────────

\subsection{行内公式与单行公式}

行内公式直接用美元符号包裹，例如 $f(x) = x^2 + 1$ 或
薛定谔方程 $\hat{H}\psi = E\psi$。

带编号的单行公式使用 \texttt{equation} 环境：
\begin{equation}
  \int_{-\infty}^{+\infty} e^{-x^2}\,\dd x = \sqrt{\pi}.
  \label{eq:gaussian}
\end{equation}

引用时用 \cref{eq:gaussian}（或 \eqref{eq:gaussian}）。

\subsection{多行对齐公式}

使用 \texttt{align} 环境在等号处对齐：
\begin{align}
  (a+b)^2 &= a^2 + 2ab + b^2, \label{eq:binom2} \\
  (a-b)^2 &= a^2 - 2ab + b^2. \label{eq:binom2m}
\end{align}

不需要编号的行在末尾加 \verb|\notag| 或改用 \texttt{align*}。

\subsection{分段函数}

\begin{equation}
  |x| =
  \begin{cases}
    \phantom{-}x, & x \geq 0, \\
    -x,           & x < 0.
  \end{cases}
\end{equation}

\subsection{矩阵}

\begin{equation}
  A = \begin{pmatrix} a_{11} & a_{12} \\ a_{21} & a_{22} \end{pmatrix},
  \qquad
  \det(A) = a_{11}a_{22} - a_{12}a_{21}.
\end{equation}

% ────────────────────────────────────────────────────────────
\section{定理、引理与证明}
% ────────────────────────────────────────────────────────────

% 定理环境在 main.tex 中已定义（theorem / lemma / proposition 等）

\begin{theorem}[勾股定理]\label{thm:pythagoras}
  设 $a$、$b$ 为直角三角形的两直角边，$c$ 为斜边，则
  \[
    a^2 + b^2 = c^2.
  \]
\end{theorem}

\begin{proof}
  （在此给出证明过程。）
\end{proof}

\begin{lemma}\label{lem:example}
  （引理的内容。）
\end{lemma}

\begin{corollary}
  由\cref{thm:pythagoras} 可知，……
\end{corollary}

% ────────────────────────────────────────────────────────────
\section{表格}
% ────────────────────────────────────────────────────────────

使用 \texttt{booktabs} 宏包（已在 \texttt{main.tex} 中加载）绘制三线表，
如\cref{tab:example}所示。

\begin{table}[htbp]
  \centering
  \caption{示例表格：不同方法在基准上的对比}
  \label{tab:example}
  \begin{tabular}{l c c c}
    \toprule
    方法           & 指标 1 & 指标 2 & 指标 3 \\
    \midrule
    方法 A         & 0.832  & 1.24   & 97.3\% \\
    方法 B         & 0.751  & 1.56   & 95.1\% \\
    \textbf{本文方法} & \textbf{0.891} & \textbf{1.10} & \textbf{98.6\%} \\
    \bottomrule
  \end{tabular}
\end{table}

% 若表格较宽，可改用 longtable 或旋转（rotating 宏包的 sidewaystable）。

% ────────────────────────────────────────────────────────────
\section{插图}
% ────────────────────────────────────────────────────────────

图片建议以 PDF 或 PNG 格式存放在 \texttt{figs/} 子目录中。
\Cref{fig:example} 展示了插图的基本用法。

\begin{figure}[htbp]
  \centering
  % 取消注释并替换文件名：
  % \includegraphics[width=0.6\textwidth]{figs/your_figure.pdf}
  \framebox[0.6\textwidth]{\rule{0pt}{4cm}\centering 图片占位}
  \caption{图片标题：说明图片内容。}
  \label{fig:example}
\end{figure}

\subsection{并排两图}

\begin{figure}[htbp]
  \centering
  % 左图
  \begin{minipage}[b]{0.45\textwidth}
    \centering
    % \includegraphics[width=\textwidth]{figs/fig_a.pdf}
    \framebox[\textwidth]{\rule{0pt}{3cm}\centering 左图}
    \caption{左图标题。}
    \label{fig:left}
  \end{minipage}
  \hfill
  % 右图
  \begin{minipage}[b]{0.45\textwidth}
    \centering
    % \includegraphics[width=\textwidth]{figs/fig_b.pdf}
    \framebox[\textwidth]{\rule{0pt}{3cm}\centering 右图}
    \caption{右图标题。}
    \label{fig:right}
  \end{minipage}
\end{figure}

% ────────────────────────────────────────────────────────────
\section{列表}
% ────────────────────────────────────────────────────────────

无序列表：
\begin{itemize}
  \item 第一条。
  \item 第二条。
  \item 第三条。
\end{itemize}

有序列表：
\begin{enumerate}
  \item 步骤一：……
  \item 步骤二：……
  \item 步骤三：……
\end{enumerate}

% ────────────────────────────────────────────────────────────
\section{文献引用与交叉引用}
% ────────────────────────────────────────────────────────────

引用文献用 \verb|\cite{key}|，例如见文献~\cite{doe2020}。
多篇同时引用：\cite{doe2020,smith2021}。

引用本文中的公式、定理、图、表分别用 \verb|\cref{}|：
\cref{eq:gaussian}、\cref{thm:pythagoras}、\cref{fig:example}、\cref{tab:example}。

% ────────────────────────────────────────────────────────────
\section{本章小结}
% ────────────────────────────────────────────────────────────

本章介绍了……（总结本章主要内容）。

\chapter{第三章标题}

\section{第一节标题}

\section{第二节标题}

\section{本章小结}

% 规范要求每章末尾设"本章小结"，概括本章主要内容及结论。

\chapter{第四章标题}

\section{第一节标题}

\section{第二节标题}

\chapter{结论与展望}

\section{主要结论}

% 总结全文的主要研究成果与创新点。

\section{研究展望}

% 指出本研究的不足之处，并对未来研究方向进行展望。
   % 结论与展望
% 如需增加章节，复制 chapt3.tex 并在此处添加 \chapter{结论与展望 / Conclusions and Outlook}

\section{结论}

\section{展望}
 等

% ---- 参考文献与附录 ----
\renewcommand{\bibname}{参考文献}
\addcontentsline{toc}{chapter}{参考文献}
{\small\setlength{\parindent}{0pt}\setlength{\parskip}{0.35em}\sloppy
\begin{thebibliography}{99}

% 在此添加参考文献条目，格式遵循 GB/T 7714 国标。
% 删除下面的示例条目，替换为自己的文献。

% 期刊论文
\bibitem{doe2020} DOE J, ROE J. An example journal article[J]. Journal of Examples, 2020, 10(2): 100--120.

% 会议论文
\bibitem{smith2021} SMITH A. An example conference paper[C]// Proceedings of the Example Conference. Beijing: Example Press, 2021: 55--63.

% 专著
\bibitem{brown2019} BROWN B. An Example Textbook[M]. New York: Example Press, 2019.

% 学位论文
\bibitem{zhang2022} 张三. 示例博士学位论文题目[D]. 北京: 北京工业大学, 2022.

\end{thebibliography}
}

\appendix
\let\cleardoublepage\clearpage

\chapter{附录A 标题}
\label{app:A}

% 在此撰写附录 A 内容。

\chapter{附录B 标题}
\label{app:B}

% 在此撰写附录 B 内容。

\cleardoublepage
\begin{publication}
\zihao{-4}\songti

% 列出在攻读学位期间发表的学术论文 / 参与的科研项目。
% 示例：
% 1. 作者甲，作者乙. 论文标题[J]. 期刊名称, 年份, 卷(期): 起止页码.

\end{publication}
   % 发表论文情况

% ---- 后置部分：致谢、索引 ----
\backmatter
\begin{acknowledgement}

% 在此撰写致谢内容。

\end{acknowledgement}


% 索引（可选；与 preamble 中的 \usepackage{makeidx} 配套）
% 正文中用 \index{词条} 标记需要收录的词条。
% 编译完成后需额外运行一次 makeindex：
%   xelatex main → makeindex main → xelatex main → xelatex main
% 如不需要索引，注释掉下面这行。
\printindex

\end{document}
